\section{Real roots of the Doche--Mend\`es France polynomials (\oeis{A347171})}\label{s:A347171}
The triangle \oeis{A347171} gives the coefficients of $P_n(y)=\sum_{0\le j<2^n}\mathrm{GRS}(j)\,y^{\mathrm{wt}(j)}$, where $\mathrm{GRS}$ is the Golay--Rudin--Shapiro sequence. The entry (K.~Ryde, 2021) records the conjecture of Doche and Mend\`es France that the number of real roots of $P_n$ is \oeis{A285869}$(n)$, the number of zeros of the Chebyshev polynomial $S_n$ in $(-\sqrt2,\sqrt2)$.

\paragraph{Prior work.} In their paper (Budapest 2001; version of April 2005, pp.~9--10), Doche and Mend\`es France
write that computations up to degree $200$ with PARI's \texttt{polsturm} suggest the count
$2\lceil n/4\rceil-((-1)^{n+1}+1)/2$, which equals \oeis{A285869}$(n)$, and add: ``Unfortunately we are not able to
establish this.'' We found no later proof.

\subsection{Statement}
The polynomials satisfy $P_0=1$, $P_1=1+y$ and
\begin{equation}\label{e:grsrec}
P_n(y)=(1-y)P_{n-1}(y)+2y\,P_{n-2}(y)\qquad(n\ge2),
\end{equation}
a recurrence due to Doche and Mend\`es France \cite{DMF}. Let
$A(n)=2b(n)+(n \bmod 2)$ with $b(n)=\lfloor n/2\rfloor-\lfloor (n+1)/4\rfloor$; this is the closed form of
\oeis{A285869}$(n)$. Put $L=-3-2\sqrt2$ and $R=-3+2\sqrt2$.

\begin{theorem}\label{t:A347171}
For every $n\ge0$ the real roots of $P_n$ are simple and their number is $A(n)$. More precisely, for
$n\ge2$, with $q=(n+3)/4$ and $K=\lceil q\rceil-1$, the number of roots of $P_n$ is
\[
\begin{array}{ll}
0 & \text{in } (-\infty,L]\ \text{and in}\ [R,0],\\
K-1+[q\in\mathbb Z] & \text{in } (L,-1),\\
{}[q\in\mathbb Z] & \text{at } y=-1,\\
K & \text{in } (-1,R),\\
{}[n \text{ even}] & \text{in } (0,\infty).
\end{array}
\]
\end{theorem}
\noindent Here $[\,\cdot\,]$ is $1$ if the condition holds and $0$ otherwise. The cases $n=0,1$ are trivial
($P_0=1$, $P_1=1+y$).

\subsection{The Chebyshev form for \texorpdfstring{$y<0$}{y<0}}
Let $S_n$ be defined by $S_{-1}=0$, $S_0=1$, $S_n(x)=xS_{n-1}(x)-S_{n-2}(x)$, so that
$S_n(2\cos\theta)=\sin((n+1)\theta)/\sin\theta$ and $S_n(2\cosh\eta)=\sinh((n+1)\eta)/\sinh\eta$.

\begin{lemma}\label{l:cheb}
For $y<0$ put $t=\sqrt{-2y}>0$ and $x=(1-y)/t=t/2+1/t$. Then $P_n(y)=t^n\bigl(S_n(x)-t\,S_{n-1}(x)\bigr)$.
\end{lemma}
\begin{proof}
Both sides equal $1$ for $n=0$ and $tx-t^2=1+y$ for $n=1$. Writing $P_n=t^nQ_n$, the recurrence \eqref{e:grsrec}
becomes $Q_n=xQ_{n-1}-Q_{n-2}$ because $tx=1-y$ and $t^2=-2y$.
\end{proof}
Since $x=t/2+1/t\ge\sqrt2$, we have $x<2$ exactly when $t\in(2-\sqrt2,2+\sqrt2)$, that is, $y\in(L,R)$.

\begin{lemma}\label{l:outside}
For $n\ge1$, $P_n$ has no root in $(-\infty,L]\cup[R,0]$.
\end{lemma}
\begin{proof}
$P_n(0)=1$. At $y=L$ or $R$ we have $x=2$, $S_n(2)=n+1$ and $P_n=t^n(n+1-tn)\ne0$ since $t=2\pm\sqrt2$
is irrational. Otherwise $x>2$; write $x=2\cosh\eta$, $\eta>0$. By Lemma~\ref{l:cheb}, $P_n(y)=0$ if and only if
$\rho:=\sinh((n+1)\eta)/\sinh(n\eta)=t$. Now $\rho=\cosh\eta+\coth(n\eta)\sinh\eta$ and
$1<\coth(n\eta)\le\coth\eta$, so $e^\eta<\rho\le2\cosh\eta=x$. If $t>2+\sqrt2$ then $\rho\le x=t/2+1/t<t$.
If $t<2-\sqrt2$ then $\rho>e^\eta>1>t$.
\end{proof}

\subsection{The phase on \texorpdfstring{$(L,R)$}{(L,R)}}
For $y\in(L,R)$ write $x=2\cos\theta$ with $\theta\in(0,\pi/4]$, and let $c=\cos\theta$, $s=\sin\theta$,
$w=\sqrt{2\cos2\theta}=\sqrt{4c^2-2}$. Since $t^2-2xt+2=0$, we get $t=2c\pm w$:
\begin{itemize}
\item branch 1: $t_1=2c+w\in(\sqrt2,2+\sqrt2)$, which describes $y\in(L,-1)$;
\item branch 2: $t_2=2c-w\in(2-\sqrt2,\sqrt2)$, which describes $y\in(-1,R)$.
\end{itemize}
On each branch $\theta\mapsto y$ is a bijection from $(0,\pi/4)$; $\theta=\pi/4$ is $y=-1$, and
$\theta\to0$ is $y\to L$ on branch 1 and $y\to R$ on branch 2. Since
$\sin((n+1)\theta)-t\sin(n\theta)=\operatorname{Im}\bigl(e^{in\theta}(e^{i\theta}-t)\bigr)$, Lemma~\ref{l:cheb} gives
\begin{equation}\label{e:phase}
P_n(y)=t^n\,\bigl|e^{i\theta}-t\bigr|\,\frac{\sin\Phi(\theta)}{\sin\theta},\qquad
\Phi(\theta)=n\theta+\beta(\theta),\quad \beta(\theta)=\arg\bigl(e^{i\theta}-t\bigr)\in(0,\pi),
\end{equation}
and $|e^{i\theta}-t|^2=1-2tc+t^2=t^2/2>0$. So the roots in $(L,R)$ are the solutions of $\Phi\in\pi\mathbb Z$.

\begin{lemma}\label{l:b2}
On branch 2, $\beta_2$ is strictly increasing from $0$ to $3\pi/4$. Hence $\Phi_2=n\theta+\beta_2$ is
strictly increasing from $0$ to $q\pi$.
\end{lemma}
\begin{proof}
$\cot\beta_2=(c-t_2)/s=(w-c)/s$. With $w'=-4sc/w$, the derivative of $(w-c)/s$ is $N/s^2$ where
$N=(w'+s)s-(w-c)c$, and $Nw=-4s^2c+w-(4c^2-2)c=w-2c<0$ because $w^2=4c^2-2<4c^2$. So $\cot\beta_2$ is
strictly decreasing. At $\theta\to0$, $\cot\beta_2\to+\infty$; at $\theta=\pi/4$, $\cot\beta_2=-1$.
\end{proof}

\begin{lemma}\label{l:b1}
On branch 1, $\beta_1'(\theta)=1-2c/w$. Hence $\beta_1$ and $\Phi_1=n\theta+\beta_1$ are strictly concave,
$\Phi_1(0^+)=\pi$, $\Phi_1(\pi/4)=q\pi$, $\Phi_1'(\theta)\to-\infty$ as $\theta\to\pi/4$, and
$\Phi_1<q\pi+\pi/4$ on $(0,\pi/4)$.
\end{lemma}
\begin{proof}
$e^{i\theta}-t_1=-(c+w)+is$ lies in the second quadrant, so $\beta_1=\pi-\arctan\bigl(s/(c+w)\bigr)$.
Differentiating,
\[
\frac{d}{d\theta}\arctan\frac{s}{c+w}=\frac{1+cw-sw'}{(c+w)^2+s^2}=\frac{w+2c}{w\,(4c^2-1+2cw)} ,
\]
using $-sw'=4s^2c/w$ and $w+cw^2+4s^2c=w+2c$. Since $(2c+w)^2=2(4c^2-1+2cw)$ and $(2c+w)(2c-w)=2$, this
equals $2/\bigl(w(2c+w)\bigr)=(2c-w)/w=2c/w-1$. Now $2c/w=2/\sqrt{4-2/c^2}$ increases strictly with $\theta$,
so $\beta_1'$ decreases strictly. The boundary values are $\beta_1(0^+)=\pi$ and $\beta_1(\pi/4)=3\pi/4$, and
$\beta_1'\to-\infty$ as $w\to0$. Finally $\theta<\pi/4$ and $\beta_1<\pi$ give $\Phi_1<n\pi/4+\pi=q\pi+\pi/4$.
\end{proof}

\subsection{Positive \texorpdfstring{$y$}{y}}
\begin{lemma}\label{l:pos}
For $y>0$, $P_n$ has exactly one root, which is simple, if $n\ge2$ is even, and no root if $n$ is odd.
\end{lemma}
\begin{proof}
Let $\rho_1>0>\rho_2$ be the roots of $f(r)=r^2-(1-y)r-2y$. Since $f(1)=-y<0<2=f(2)$, $1<\rho_1<2$. The
recurrence \eqref{e:grsrec} at $n=1$ gives $P_{-1}=1$, hence
$P_n=\bigl(\rho_1^{n+1}(\rho_2-1)+\rho_2^{n+1}(1-\rho_1)\bigr)/(\rho_2-\rho_1)$, and $P_n(y)=0$ if and only if
$\rho_1^{n+1}(1-\rho_2)=\rho_2^{n+1}(1-\rho_1)$. The left side is positive. If $n$ is odd the right side is
negative. If $n$ is even, parametrize by $\rho=\rho_1\in(1,2)$: then $y=\rho(\rho-1)/(2-\rho)$ is an increasing
bijection onto $(0,\infty)$, $|\rho_2|=2(\rho-1)/(2-\rho)$, and the equation becomes
$\kappa^{n+2}=2/(2-\rho)^2$ with $\kappa=2(\rho-1)/(\rho(2-\rho))$. Let
$g(\rho)=(n+2)\log\kappa+2\log(2-\rho)-\log2$. Then
$g'(\rho)=(n+2)\bigl(\tfrac1{\rho-1}-\tfrac1\rho\bigr)+\tfrac{n}{2-\rho}>0$, $g\to-\infty$ at $1^+$ and, for
$n\ge1$, $g\to+\infty$ at $2^-$. So there is exactly one root, and it is simple.
\end{proof}

\subsection{Proof of Theorem \ref{t:A347171}}
Let $n\ge2$, so $q>1$. By \eqref{e:phase} and Lemma~\ref{l:b2}, the roots in $(-1,R)$ correspond to the integers
$k$ with $0<k<q$: there are $K$ of them. At $y=-1$, $\Phi=q\pi$, so $-1$ is a root if and only if
$q\in\mathbb Z$, that is, $n\equiv1\pmod4$. On branch 1, by Lemma~\ref{l:b1}, $\Phi_1$ is strictly concave with
boundary values $\pi<q\pi$, its maximum $M$ is interior with $q\pi<M<q\pi+\pi/4$, and the only multiple of
$\pi$ in $[q\pi,M]$ is $q\pi$ itself when $q\in\mathbb Z$ (otherwise the next multiple is at least
$(q+1/4)\pi$). By concavity, each $k\pi$ with $1<k<q$ is attained exactly once, on the increasing part; if
$q\in\mathbb Z$, $q\pi$ is attained once in the interior and once at the endpoint $\theta=\pi/4$; and $\pi$ is
not attained in the interior. This gives $K-1+[q\in\mathbb Z]$ roots in $(L,-1)$. All these roots are simple,
because $\Phi'\ne0$ there (at $y=-1$ one uses the parameter $t$, in which $\theta'(\sqrt2)=0$ and $\beta'\ne0$).
With Lemmas~\ref{l:outside} and~\ref{l:pos}, the total is $2K-1+2[q\in\mathbb Z]+[n\text{ even}]$.
Writing $n=4m+r$: for $r=0,1,2,3$ we get $(q,K)=(m+\tfrac34,m)$, $(m+1,m)$, $(m+\tfrac54,m+1)$,
$(m+\tfrac32,m+1)$ and totals $2m$, $2m+1$, $2m+2$, $2m+1$, while $b(n)=m,m,m+1,m$. In each case the total
is $A(n)$. \qed


\begin{remark}
Every formula above was checked numerically, and the table of Theorem~\ref{t:A347171} was confirmed for
$1\le n\le300$ with exact Sturm sequences, the signs at $L$ and $R$ being computed in $\mathbb Q(\sqrt2)$.
The scripts are in the repository \url{https://github.com/joaquimmj/oeis-walnut}, folder
\texttt{conjectures/A347171}.
\end{remark}

